\subsection{Milestone 4: Model Deployment -- API}
\label{sec:m4}
\guidance{
  \textbf{Template:} Describe how you have deployed your model and the technologies used.
}

\subsubsection{ML-based component/system architecture}
\label{sec:m4-architecture}
\guidance{
  \textbf{Template:} Physical and logical architectures of the system, including diagrams.
  Design patterns used.
  Selection and justification of design decisions (e.g.\ cloud provider).
  The demo repo compares free tiers of AWS, Azure, GCP, Okteto, and the FIB Virtech VM (free for the semester, 4\,GB RAM, 20\,GB storage); see references/course-demos.md.\\
  \textbf{Rubric (15 pts, the largest single item):} In what ways does our system design balance clarity and justified complexity while applying suitable architectural patterns and principles?\\
  Component responsibilities and interactions are typical M4 EDN entries.
}

\tbd{logical and physical architecture with diagrams}

\paragraph{Cloud provider.}
\label{sec:m4-cloud}
We deploy on the FIB Virtech virtual machine that the course provides (EDN-17).
The course lists four free alternatives: AWS, Azure, GCP, and Okteto.

Memory decided the choice, together with the fact that the course provides the machine.
Our stack is an API container, Prometheus, Grafana, and a scheduled drift job, which does not fit into the 1\,GB and 0.6\,GB of the AWS and GCP free tiers.
Okteto offers 3\,GB and Azure between 1 and 4\,GB, depending on the instance, so neither has more memory than the 4\,GB of the FIB machine.
The FIB machine also has 20\,GB of disk and is free for the semester the project runs.
Because the course provides it, nobody needs a personal cloud account with payment details (EDN-17).

NFR-04 gives the machine a budget: the API container stays below 1\,GB of memory, the drift job runs in its own container, and the metrics and logs are capped so that the disk cannot fill up.
NFR-12 keeps the choice reversible: the system starts on a fresh machine with one command and needs no paid service, so nothing in the serving path depends on this provider.

The machine is a single host without redundancy, which is why NFR-05 asks for fast recovery and not for an uptime percentage.
Nobody on the team has access to it yet; if this is still the case at the first Milestone~4 laboratory session, we raise it with the lecturers.

\subsubsection{API design}
\label{sec:m4-api-design}
\guidance{
  \textbf{Template:} e.g.\ FastAPI.\\
  \textbf{Rubric (10 pts, shared with API testing):} To what extent do our API design and its documentation make the system robust, intuitive, and easy to adopt by others?\\
  The demo separates the FastAPI app from its Pydantic request/response schemas and relies on the generated Swagger UI (\texttt{/docs}) and ReDoc (\texttt{/redoc}).
}

\tbd{endpoints, schemas and API documentation}

\subsubsection{API testing}
\label{sec:m4-api-testing}
\guidance{
  \textbf{Template:} e.g.\ FastAPI, pytest.\\
  The demo tests the endpoints with FastAPI's \texttt{TestClient}, covering the happy path and invalid input.
}

\tbd{API tests}
